\section{Point models and approximate intertwining}\label{pt:section}

We now construct homomorphisms \(P\to Q\) and \(Q\to P\),
and then an isomorphism realizing the prescribed invariant.
Cone existence first gives a model on the line. Norm uniqueness
makes translation covariant, so the model extends to a crossed
product. A trace-normalized projection then gives a point model.
A second application of uniqueness makes its late coordinate
lifts close modulo unitaries, yielding an actual homomorphism.
The final application makes the two compositions approximate
identities and completes the intertwining.

The suspension, crossed-product projection, and subsequence
extraction also occur in \cite[Remark~6.4]{GabeO2}. The suspension,
crossed-product, and full-corner strategy adapts the proof of
\cite[Theorem~14.1]{Gabe}. The latter theorem has an
\(\mathcal O_\infty\)-stable target and is not applied directly here.
Here each passage must preserve the
full extended trace data. We compute the projection's trace by
an identity of finite positive sums, and recover the invariant
of the limiting maps through positive cuts.
We use \(\Ad(w)(a)=waw^*\).

\subsection{The line and its translation}

\begin{proposition}[A covariant line model]\label{pt:suspension}
There is a model
\[
 \lambda:C_0(\R)\otimes P\longrightarrow D
\]
for Lebesgue measure and \(F\).
Let \(\alpha(f)(t)=f(t-1)\), acting trivially on \(P\).
For c.p.c. coordinate lifts \(\lambda_n\), there are
unitaries \(u_n\in M(Q)\) such that
\begin{equation}\label{pt:coordinate-covariance}
 \|\Ad(u_n)(\lambda_n(s))-\lambda_n(\alpha(s))\|
 \longrightarrow0\qquad(s\in C_0(\R)\otimes P).
\end{equation}
\end{proposition}

\begin{proof}
For \(j\geq2\), put \(I_j=(-j,j]\) and give it Lebesgue
measure.  This is a finite full-support cone measure, so
\Cref{cc:cone-existence}, transported along a homeomorphism
\((0,1]\cong I_j\), gives a model
\(\lambda^{(j)}:C_0(I_j)\otimes P\to D\).
Choose \(0\leq\eta_j\leq1\) in \(C_c((-j,j))\), equal
to one on \([-j+1,j-1]\).  The maps
\[
 R_j:C_0(\R)\otimes P\longrightarrow C_0(I_j)\otimes P,
 \qquad R_j(s)=\eta_j\,s|_{I_j},
\]
are c.p.c.: restriction followed by multiplication by the scalar
positive contraction $\eta_j$ is completely positive and contractive,
and its values vanish at the omitted left endpoint. The domain
$C_0(I_j)\otimes P$ of each $\lambda^{(j)}$ is separable and nuclear.
The maps $R_j$ are exactly multiplicative on every
fixed finite set of compactly supported inputs once \(j\)
is large enough to contain their supports in the region
where \(\eta_j=1\).

Take c.p.c. coordinate lifts of each \(\lambda^{(j)}\),
by \cite[Theorem~3.10]{ChoiEffrosLifting}, and compose with \(R_j\).
Apply \Cref{seq:diagonal} to an increasing dense list of
compactly supported algebra tests and the elementary rank
tests of \Cref{seq:elementary-tests}.
Here the tests are fixed on the line before increasing
\(j\): each function, its positive trims, and its enlarged
compact support all lie inside \((-j+1,j-1)\) in every
sufficiently late row.  Their Lebesgue masses are therefore
unchanged by \(R_j\).
For every sandwich test choose the strict scalar slack,
the enlarged source, and the trimmed target before selecting
coordinates in that row.  The constant Cuntz representatives
in \(Q\) are the same for all these rows.
The fixed target tolerances give the bounded comparison
witnesses of \Cref{seq:uniform-witness}; thus the row
diagonal retains the comparisons in \(D\), even when the
row thresholds vary.
Choose the coordinates for the first \(j\) tests only
after these choices, with algebra errors at most \(j^{-1}\).
The resulting c.p.c. sequence \(\lambda_n\) is
asymptotically multiplicative.  Its quotient
\(\lambda\) satisfies the elementary rank identities for
every regularized limit, and
\Cref{seq:elementary-tests} proves the asserted model
identity.

Lebesgue measure is translation invariant, including the
extended product convention.  Hence \(\lambda\) and
\(\lambda\circ\alpha\) are models for the same data.
Apply \Cref{un:uniqueness} to their specified coordinate
lifts to obtain \eqref{pt:coordinate-covariance}.
\end{proof}

The covariance unitaries belong to \(M(Q)\), whereas the coordinate model
values belong to \(Q\). To place them in one unital algebra, let
\[
 \mathcal C=\prod_{n=1}^{\infty}M(Q)\Big/
             \bigoplus_{n=1}^{\infty}M(Q).
\]
The product here consists of bounded sequences and the direct sum
of norm-null sequences. There is a natural faithful inclusion
\(D\subseteq\mathcal C\), and \(D\) is an ideal in
\(\mathcal C\).
The sequence from \Cref{pt:suspension} defines a unitary
\(u\in\mathcal C\), and the covariance relation is exact:
\(\Ad(u)\lambda(s)=\lambda(\alpha(s))\).
Integration of this covariant pair gives a homomorphism
\begin{equation}\label{pt:crossed-map}
 \Psi:(C_0(\R)\otimes P)\rtimes_\alpha\Z
       \longrightarrow D.
\end{equation}
This also applies when \(\lambda\) is degenerate: in a faithful
representation of \(\mathcal C\) on \(\mathcal H\), restrict to the reducing
subspace \(\overline{\lambda(C_0(\R)\otimes P)\mathcal H}\),
integrate there, and extend by zero. The finite Fourier sums
have the form \(\sum_n\lambda(s_n)u^n\), which lie in the
ideal \(D\); norm closure gives the stated target.
Because the action on \(P\) is trivial and \(\Z\) is
amenable, the domain is canonically
\[
 E\otimes P,\qquad E=C_0(\R)\rtimes_\alpha\Z.
\]
The identification follows also by matching the commuting
representations of \(P\) and the covariant representation
of \(C_0(\R)\).  Nuclearity makes the tensor-product norm
unambiguous.  We denote the canonical implementing
multiplier of \(E\) by \(v\), so \(vf v^*=\alpha(f)\).

\subsection{An explicit trace-normalized corner}

For any projection \(p\in E\), the map \(a\mapsto p\otimes a\)
is a homomorphism from \(P\) to \(E\otimes P\). We seek a
projection for which composition with \(\Psi\) has exactly the
point trace data. The next lemma supplies the needed normalization
by a finite positive partition identity. It also identifies the
corresponding full corner of \(E\).

\begin{lemma}[The corner projection]\label{pt:corner}
There are \(\zeta\in C_c(\R)_+\) and a projection \(p\in E\)
such that
\[
 \sum_{n\in\Z}\zeta(t-n)^2=1,
 \qquad \int_\R\zeta(t)^2\,dt=1,
\]
and, for any \(C^*\)-algebra \(R_0\), any extended trace
\(\omega\) on \(E\otimes R_0\), and every \(a\in(R_0)_+\),
\begin{equation}\label{pt:corner-trace}
 \omega(p\otimes a)=\omega(\zeta^2\otimes a).
\end{equation}
In particular, if the restriction of \(\omega\) to
\(C_0(\R)\otimes R_0\) is Lebesgue measure tensored with
\(\nu\), then \(\omega(p\otimes a)=\nu(a)\), including
the value infinity. Moreover, \(p\) is full and
\(pEp\cong C(\R/\Z)\).
\end{lemma}

\begin{proof}
Choose \(\eta\in C_c(\R)_+\) strictly positive on
\([0,1]\), and set
\[
 H(t)=\sum_{n\in\Z}\eta(t-n)^2,\qquad
 \zeta(t)=H(t)^{-1/2}\eta(t).
\]
The sum defining \(H\) is locally finite, positive,
continuous, and one-periodic.  The partition identity
for \(\zeta\) follows immediately.  Integrating that
identity over \([0,1]\) gives the asserted normalization.
Put
\begin{equation}\label{pt:projection-formula}
 p=\sum_{n\in\Z}p_n v^n,
 \qquad p_n(t)=\zeta(t)\zeta(t-n).
\end{equation}
Only finitely many coefficients are nonzero.  The
involution formula gives
\(\alpha_n(p_{-n})=p_n\), so \(p^*=p\).
Convolution gives, for every \(k\in\Z\),
\[
 (p^2)_k(t)
 =\zeta(t)\zeta(t-k)\sum_{n\in\Z}\zeta(t-n)^2
 =p_k(t).
\]
Thus \(p\) is a projection.

Choose finitely many real functions \(q_i\in C_c(\R)\),
each supported in an interval of diameter less than one,
such that
\[
 q:=\sum_iq_i^2=1\quad\hbox{on }\supp(\zeta).
\]
For completeness, choose bump functions \(r_i\) in such
intervals whose squared sum is positive near
\(\supp(\zeta)\), and a cutoff \(\kappa\) equal to one
on that support and supported where this sum is positive.
Then \(q_i=\kappa r_i/(\sum_jr_j^2)^{1/2}\), extended
by zero, have the required properties.
The coefficient formula for \(p\) yields
\begin{equation}\label{pt:positive-partition}
 qp=p=pq,\qquad \sum_iq_i p q_i=\zeta^2.
\end{equation}
To verify the second identity, its \(n\)-th coefficient
is
\[
 \zeta(t)\zeta(t-n)\sum_iq_i(t)q_i(t-n).
\]
This vanishes when \(n\ne0\), because each support has
diameter less than one, and equals \(\zeta(t)^2\) when
\(n=0\).

For \(a\in(R_0)_+\), set \(y_i=q_i p\otimes a^{1/2}\).
These are elements of \(E\otimes R_0\), and
\eqref{pt:positive-partition} gives the exact positive
identities
\[
 \sum_i y_i^*y_i=p\otimes a,\qquad
 \sum_i y_i y_i^*=\zeta^2\otimes a.
\]
Additivity and \(\omega(y_i^*y_i)=\omega(y_i y_i^*)\)
prove \eqref{pt:corner-trace}.  There are only finitely
many positive summands.  In particular, this calculation
uses no subtraction of infinite values and requires no
finiteness assumption on \(\omega\).
For the Lebesgue trace assertion, the extended product formula
on the right has value
\((\int\zeta^2)\nu(a)=\nu(a)\).

We also verify fullness and the corner description.
Equation~\eqref{pt:positive-partition} puts
\(\zeta^2\) in \(\Ideal_E(p)\).
Conjugation by \(v^n\) preserves this ideal, so it contains
every \(\alpha_n(\zeta^2)\).
For \(f\in C_c(\R)\), the partition identity expresses
\[
 f=\sum_{n\in\Z}f\,\alpha_n(\zeta^2)
\]
as a finite sum of elements of that ideal.
It therefore contains \(C_0(\R)\), and then the dense
finite Fourier sums in \(E\).  Hence \(p\) is full.

For \(f\in C_c(\R)\) and \(n\in\Z\), direct convolution
gives
\begin{equation}\label{pt:corner-compression}
 p(fv^n)p=g_{f,n}p,\qquad
 g_{f,n}(t)=\sum_{k\in\Z}
       \zeta(t-k)f(t-k)\zeta(t-k-n).
\end{equation}
The function \(g_{f,n}\) is continuous and one-periodic.
Conversely, for a continuous one-periodic \(g\), the
element \(gp\) belongs to \(E\), since all its Fourier
coefficients are compactly supported. Periodicity makes \(g\)
a central multiplier of \(E\), so \(gp\in pEp\), and the map
\(g\mapsto gp\) is a unital homomorphism
\(C(\R/\Z)\to pEp\).  It is injective: its zeroth
coefficient is \(g\zeta^2\), and every orbit has a
point where \(\zeta\ne0\).
Equation~\eqref{pt:corner-compression} and density of
compactly supported Fourier sums prove surjectivity.
\end{proof}

\Needspace{3\baselineskip}
\begin{corollary}[Existence of point models]\label{pt:point-model}
There is a homomorphism \(\varphi:P\to Q_\infty\) such that
\[
 \rho\circ\varphi=F(\rho|_Q)
\]
for every regularized limit \(\rho\).
\end{corollary}

\begin{proof}
Use \eqref{pt:crossed-map} and \Cref{pt:corner} to define
\(\varphi(a)=\Psi(p\otimes a)\).
For a regularized limit \(\rho\), the extended trace
\(\omega=\rho\circ\Psi\) restricts on
\(C_0(\R)\otimes P\) to Lebesgue measure tensored with
\(F(\sigma)\), where \(\sigma=\rho|_Q\).
Thus, for every \(a\in P_+\),
\[
 \rho(\varphi(a))
 =\omega(p\otimes a)
 =\omega(\zeta^2\otimes a)
 =F(\sigma)(a).
\]
The positive partition calculation applies equally when
this common value is infinite.
\end{proof}

\subsection{From a point model to an actual map}

The coordinate lifts of a point model need not converge.
Uniqueness will make their unitary orbits Cauchy on finite sets.
We can then conjugate a subsequence so that its successive
errors are summable and its limit lies in \(Q\).

\begin{proposition}[An actual homomorphism]\label{pt:actual-map}
There is a homomorphism \(\phi:P\to Q\) satisfying
\[
 \tau\circ\phi=F(\tau)\qquad(\tau\in T(Q)),
 \qquad \Cu(\phi)=G.
\]
It is injective, and its constant sequence is a point model.
\end{proposition}

\begin{proof}
Take c.p.c. lifts \(\varphi_n:P\to Q\) of the point
model in \Cref{pt:point-model}, using nuclearity and
\cite[Theorem~3.10]{ChoiEffrosLifting}.
For every finite \(\mathcal E\subseteq P\) and every
\(\eps>0\), there is \(N\) such that
\begin{equation}\label{pt:all-pairs}
 \inf_{w\in\mathcal U(M(Q))}
 \max_{a\in\mathcal E}
 \|\Ad(w)(\varphi_m(a))-\varphi_n(a)\|<\eps
 \qquad(m,n\geq N).
\end{equation}
If this failed, choose pairs \(m_k,n_k\to\infty\),
with both index sequences strictly increasing,
for which the displayed infimum is at least \(\eps\).
Every subsequence of a model is a model for the same
data: given coordinate traces $\tau_k$ on indices $n_k$ and a free
ultrafilter $\omega$ on $k$, set $\widetilde\tau_{n_k}=\tau_k$
and use the zero trace at the remaining coordinates. Define
$A\in\widetilde\omega$ if and only if
$\{k:n_k\in A\}\in\omega$. Then $\widetilde\omega$ is a free
ultrafilter concentrated on the chosen indices, and every cut
ultralimit agrees with the corresponding subsequence ultralimit.  Thus \((\varphi_{m_k})\) and
\((\varphi_{n_k})\) define two point models.
\Cref{un:uniqueness} contradicts the choice of the pairs.

Let \(\mathcal E_k\) be increasing finite sets with dense
union in the unit ball of \(P\).
Use \eqref{pt:all-pairs} to choose an increasing
subsequence \(n_k\) and \(w_k\in\mathcal U(M(Q))\) with
\[
 \max_{a\in\mathcal E_k}
 \|\Ad(w_k)(\varphi_{n_{k+1}}(a))
                -\varphi_{n_k}(a)\|<2^{-k}.
\]
Choose \(n_k\) beyond the threshold for the \(k\)-th
test, so that both indices in this inequality are
eligible.  Starting with \(W_1=1\), set
\(W_{k+1}=W_k w_k\) and
\(\psi_k=\Ad(W_k)\circ\varphi_{n_k}\).
The sequence \(\psi_k(a)\) is Cauchy for every \(a\)
in the dense union, with summable successive errors,
and then for every \(a\in P\), by contractivity.
Its limit \(\phi:P\to Q\) is c.p.c. and multiplicative,
since the multiplicative defects of the original
coordinate lifts tend to zero.  It is therefore a
homomorphism.

We check the invariant without assuming continuity of
unbounded trace evaluation under norm convergence.
Fix \(\tau\in T(Q)\), and let \(\rho_\tau\) be its
regularized constant-coordinate limit.
The maps \(\psi_k\) and the constant map \(\phi\)
give the same homomorphism into \(Q_\infty\).
For every \(a\in P_+\) and every cut tolerance,
coordinate unitary invariance gives
\[
 \tau((\psi_k(a)-\eps)_+)
   =\tau((\varphi_{n_k}(a)-\eps)_+).
\]
Regularization and the subsequence model identity
therefore imply
\[
 \tau(\phi(a))
 =\rho_\tau((\psi_k(a))_k)
 =F(\tau)(a).
\]
Here \(\rho_\tau|_Q=\tau\) follows from
\Cref{seq:regularized}.  Applying the equality to
positive powers gives
\[
 d_\tau(\phi(a))=d_{F(\tau)}(a)
               =d_\tau(G[a]).
\]
This calculation also applies to finite matrices. If
\(x=(x_{ij})\in M_k(P)_+\), the unnormalized matrix traces give
\begin{align*}
 (\tau\otimes\Tr_k)(\phi^{(k)}(x))
 &=\sum_i\tau(\phi(x_{ii}))
 =\sum_i F(\tau)(x_{ii})\\
 &=(F(\tau)\otimes\Tr_k)(x),
\end{align*}
including infinite values. Apply this identity to positive
powers and then use positive cuts of finite-matrix
approximations in \(P\otimes\K\). Pointwise trace
determination in \(\Cu(Q)\) gives \(\Cu(\phi)=G\).
Since \(G\) is injective, \(\phi(a)=0\) for positive
\(a\) implies \([a]=0\), hence \(a=0\);
thus \(\phi\) is injective.
Finally, for any regularized limit \(\rho\), put
\(\sigma=\rho|_Q\).  Its value on the constant map is
\[
 \rho\circ\phi=\sigma\circ\phi=F(\sigma).
\]
So the constant map is itself a point model.
\end{proof}

\subsection{The final intertwining}

\begin{theorem}[Trace-cone classification]\label{pt:intertwining}
There is an isomorphism \(\Phi:P\to Q\) such that
\[
 \tau\circ\Phi=F(\tau)\quad(\tau\in T(Q)),
 \qquad \Cu(\Phi)=G.
\]
Consequently \Cref{thm:main} holds.
\end{theorem}

\begin{proof}
Apply \Cref{pt:actual-map} to \(F\) and to \(F^{-1}\).
This gives homomorphisms
\(\phi:P\to Q\) and \(\psi:Q\to P\) with inverse
trace and Cuntz data.
Their compositions have identity trace data.  The
constant maps \(\psi\phi\) and \(\id_P\) are therefore
point models for the same data, and similarly for
\(\phi\psi\) and \(\id_Q\).
By \Cref{un:uniqueness}, the compositions can be
conjugated arbitrarily close to the corresponding
identity on each finite set.
We give the intertwining construction to ensure that
the conclusion is an isomorphism in \(P,Q\).

Fix increasing finite sets \(\mathcal P_n,\mathcal Q_n\)
with dense unions in the unit balls, and put
\(\eps_n=2^{-n}\).
Start with \(f_1=\phi\), \(g_0=\psi\), and empty
\(\mathcal E_0,\mathcal F_0\).
Suppose \(f_n\) and \(g_{n-1}\) have been chosen.
Choose a finite set
\[
 \mathcal E_n\supseteq
 \mathcal E_{n-1}\cup\mathcal P_n
       \cup g_{n-1}(\mathcal F_{n-1}).
\]
The composition \(g_{n-1}f_n\) still has identity
trace data, because every adjustment made below is
unitary conjugacy.
Choose \(v_n\in\mathcal U(M(P))\), and set
\(g_n=\Ad(v_n)\circ g_{n-1}\), so that
\begin{equation}\label{pt:left-intertwining}
 \|g_nf_n(a)-a\|<\eps_n
 \qquad(a\in\mathcal E_n).
\end{equation}
Next choose a finite set
\[
 \mathcal F_n\supseteq
 \mathcal F_{n-1}\cup\mathcal Q_n\cup f_n(\mathcal E_n).
\]
Choose \(z_n\in\mathcal U(M(Q))\), and set
\(f_{n+1}=\Ad(z_n)\circ f_n\), with
\begin{equation}\label{pt:right-intertwining}
 \|f_{n+1}g_n(b)-b\|<\eps_n
 \qquad(b\in\mathcal F_n).
\end{equation}
This second adjustment is possible by the same
identity-data argument for \(f_ng_n\).

The inclusions of the image tests make these sequences
Cauchy.  If \(a\in\mathcal E_n\), contractivity and
\eqref{pt:left-intertwining}--\eqref{pt:right-intertwining}
give
\begin{align*}
 \|f_{n+1}(a)-f_n(a)\|
 &\leq\|f_{n+1}(a-g_nf_n(a))\|\\
 &\quad+\|f_{n+1}g_n(f_n(a))-f_n(a)\|
 <2\eps_n.
\end{align*}
If \(n\geq2\) and \(b\in\mathcal F_{n-1}\), then
\(g_{n-1}(b)\in\mathcal E_n\), and
\begin{align*}
 \|g_n(b)-g_{n-1}(b)\|
 &\leq\|g_n(b-f_ng_{n-1}(b))\|\\
 &\quad+\|g_nf_n(g_{n-1}(b))-g_{n-1}(b)\|
 <\eps_{n-1}+\eps_n.
\end{align*}
These summable bounds and density give point-norm
limits \(f_\infty:P\to Q\) and
\(g_\infty:Q\to P\), both homomorphisms.
Equations~\eqref{pt:left-intertwining} and
\eqref{pt:right-intertwining} imply
\[
 g_\infty f_\infty=\id_P,\qquad
 f_\infty g_\infty=\id_Q.
\]
For example, contractivity controls the moving input
\(f_n(a)\) by its limit \(f_\infty(a)\), after which
point-norm convergence of \(g_n\) applies to the fixed
limit input.  Density gives these equalities on the
whole algebras. Thus \(f_\infty\) is the required isomorphism.

The isomorphism realizes the prescribed trace data as well. Each
\(f_n\) has Cuntz data \(G\), since the changes to \(\phi\) are
unitary conjugacies. For \(a\in P_+\) and \(\eps>0\), choose \(n\)
with \(\|f_n(a)-f_\infty(a)\|<\eps\). Positive cut comparison gives,
for every \(\tau\in T(Q)\),
\[
 d_{F\tau}((a-\eps)_+)\leq d_\tau(f_\infty(a)),\qquad
 d_\tau((f_\infty(a)-\eps)_+)\leq d_{F\tau}(a).
\]
Taking the supremum over cuts proves equality of ranks; spectral
integration gives \(\tau\circ f_\infty=F(\tau)\), including infinite
values. This argument uses no continuity of an unbounded weight on
an uncut norm-convergent sequence. The matrix argument in
\Cref{pt:actual-map} also gives \(\Cu(f_\infty)=G\).
Set \(\Phi=f_\infty\). Substituting
\(P=A\otimes\W\otimes\K\) and
\(Q=B\otimes\W\otimes\K\) proves \Cref{thm:main}.
\end{proof}
